\section{Nearest-work challenge, challenger archives, and invention readiness}

Repeated failure is not read as an instruction to invent a new operator. Before structural invention, ordinary causes, known mechanisms, transfer candidates and mature parent-domain treatments are searched for first. The invention-readiness gate asks whether relevant search axes are bounded-flat, ordinary causes have been excluded to the supported extent, the residual survives distinct variations, and discriminator evidence identifies a missing-structure class rather than a generic score deficit.

Candidate method changes are retained as challengers rather than silently overwriting the incumbent. The archive records candidate identity, generating issue/evidence, attempted intervention, execution result, regression result, fresh-transfer result, protected-assurance outcome and final host disposition. Failed and harmful candidates remain queryable.

This structure absorbs useful mechanisms from evolutionary and self-improving-agent work without inheriting self-authorization. Archives, iterative candidate generation, population search and evaluator-guided optimization may all be valuable proposal mechanisms. They become method authority only through the separate assurance path.

\subsection{Method-level challengers and method-space expansion}
A method challenger is a candidate structural change linked to documented failure or obstruction records, not a claim that a new method has been invented. The versioned method-challenger record binds the subject revision, generating failures, competing ordinary causes already challenged, known-method search routes and their inadequacy, assimilated donor mechanisms, the structural edit and its predeclared discriminator. Candidate generation is therefore separate from host adoption. A candidate must traverse
\[
\begin{aligned}
\textsc{Diagnose}&\rightarrow\textsc{Search/Absorb}\rightarrow\textsc{Challenger}\rightarrow\textsc{IsolatedRun}\\
&\rightarrow\textsc{Replay}\rightarrow\textsc{Fresh}\rightarrow\textsc{Protected}\rightarrow\textsc{HostDisposition}.
\end{aligned}
\]
A replay improvement cannot compensate for harmful fresh transfer, and missing protected evidence blocks disposition. Rejected, null and harmful histories remain append-only and queryable; neither a generator nor a challenger can rewrite the evaluator or authority registry. For example, a structural adaptation that improves its motivating replay but harms a protected family on fresh transfer remains rejected, even if its development score is higher. Known-method search and transfer must be challenged before an invention gap is declared; repeated failure alone never makes a method ready for invention. The host alone may adopt or recommend adoption.

This is a governance boundary, not an invention result. The mechanism can govern whether a method candidate is safely evaluated and presented to a host; whether a learned system can generate valuable new method structures is a different question, and nothing here bears on it.

\subsection{Transferred design knowledge}

A prior specification audit re-examined the mature method surfaces this design inherits, treating them as prior design knowledge rather than as results. Fourteen provisional structural dimensions across 59 registered mechanics produced 826 explicit transfer questions; evidence-bound answer records moved that registered structural frontier from 826 to 0 with attributable answer and audit closure.

This result is deliberately narrow:
\[
\text{specification closure}\neq\text{empirical validity}\neq\text{promotion authority}.
\]
It can reduce rediscovery of already specified mechanics and expose remaining empirical work, but it cannot establish that the transferred contracts improve behaviour on fresh tasks.
